se comprobó experimentalmente la validez y precisión del conversor digital a analógico de 3 bits, logrando transformar las ocho combinaciones binarias de entrada en niveles de tensión analógica correspondientes. La respuesta medida en las configuraciones de solo una fuente (001, 010 y 100) coincidió con el teorico mediante el teorema del equivalente de Thévenin, mientras que las lecturas en las combinaciones de estas validaron la aplicabilidad del principio de superposición, registrando un error relativo porcentual máximo de apenas 0.40\%. Finalmente, las pequeñas desviaciones respecto a los valores enteros teóricos se justifican por la tolerancia nominal del 5\% en las resistencias de 1 \unit{\kilo\ohm}, pequeñas inestabilidades al ajustar la fuente DC y la resolución de las escalas del multímetro digital.
